\subsection{拓扑群中的换位子}

设 \(G\) 是拓扑群。我们按下列公式定义群 \(\overline{D^0}G, \overline{D^1}G, \overline{D^2}G, \ldots\) 和 \(\overline{C^1}G, \overline{C^2}G, \overline{C^3}G, \ldots\)：

\[
\begin{array}{l l} \overline {{\mathrm{D} ^ {0}}} \mathrm{G} = \mathrm{G}, & \overline {{\mathrm{D} ^ {i + 1}}} \mathrm{G} = (\overline {{\overline {{\mathrm{D} ^ {i}}} \mathrm{G}}}, \overline {{\mathrm{D} ^ {i}}} \mathrm{G}) \\ \overline {{\mathrm{C} ^ {1}}} \mathrm{G} = \mathrm{G}, & \overline {{\mathrm{C} ^ {i + 1}}} \mathrm{G} = (\overline {{\mathrm{G}}}, \overline {{\mathrm{C} ^ {i}}} \mathrm{G}). \end{array}
\]

命题 1. 设 G 是拓扑群，A 和 B 是 G 的子群。则 \((\overline{\mathrm{A}},\overline{\mathrm{B}}) = \overline{(\overline{\mathrm{A}},\overline{\mathrm{B}})},\overline{\mathrm{D}^{i}}\overline{\mathrm{A}} = \overline{\mathrm{D}^{i}\mathrm{A}},\overline{\mathrm{C}^{i}}\overline{\mathrm{A}} = \overline{\mathrm{C}^{i}\mathrm{A}}\) 。

令 \(\phi\) 为从 \((x,y)\mapsto x^{-1}y^{-1}xy\) 到 \(\mathbf{G}\times \mathbf{G}\) 的连续映射 \(\mathbf{G}\)。则 \(\phi (\mathbf{A}\times \mathbf{B})\subset (\mathbf{A},\mathbf{B})\)，从而 \(\phi (\overline{\mathbf{A}}\times \overline{\mathbf{B}})\subset (\overline{\mathbf{A}},\overline{\mathbf{B}})\)，因而

\[
\overline{(\overline{\mathrm{A}},\overline{\mathrm{B}})} \subset (\overline{\mathrm{A}},\overline{\mathrm{B}});
\]

反向包含显然成立，因此 \((\overline{\mathrm{A}},\overline{\mathrm{B}}) = \overline{(\overline{\mathrm{A}},\overline{\mathrm{B}})}\)。显然 \(\overline{\mathrm{D}^{0}}\overline{\mathrm{A}} = \overline{\mathrm{D}^{0}\mathrm{A}}\)；假设等式 \(\overline{\mathrm{D}^{i}}\overline{\mathrm{A}} = \overline{\mathrm{D}^{i}\mathrm{A}}\) 成立，则

\[
\overline{\mathrm{D}^{i+1}}\overline{\mathrm{A}} = (\overline{\overline{\mathrm{D}^{i}}\overline{\mathrm{A}}},\overline{\mathrm{D}^{i}}\overline{\mathrm{A}}) = (\overline{\mathrm{D}^{i}}\overline{\mathrm{A}},\overline{\mathrm{D}^{i}}\overline{\mathrm{A}}) = \overline{(\mathrm{D}^{i}\mathrm{A},\mathrm{D}^{i}\mathrm{A})} = \overline{\mathrm{D}^{i+1}\mathrm{A}}
\]

从而对于所有 \(\overline{\mathrm{D}^{i}}\overline{\mathrm{A}} = \overline{\mathrm{D}^{i}\mathrm{A}}\)，都有 \(i\)。等式 \(\overline{\mathrm{C}^{i}}\overline{\mathrm{A}} = \overline{\mathrm{C}^{i}\mathrm{A}}\) 的证明类似。

推论 1. 若 G 是 Hausdorff 群，则下列条件等价：

\begin{enumerate}
\def\labelenumi{(\roman{enumi})}
\item
  G 是可解的（分别地，幂零的）；
\item
  当 i 充分大时，\(\overline{\mathrm{D}^{i}}\mathrm{G} = \{e\}\)（分别地，\(\overline{\mathrm{C}^{i}}\mathrm{G} = \{e\}\)\(i\)）。
\end{enumerate}

\(\mathrm{D}^{i}\mathrm{G} \subset \overline{\mathrm{D}^{i}}\mathrm{G}, \mathrm{C}^{i}\mathrm{G} \subset \overline{\mathrm{C}^{i}}\mathrm{G}\)、，因而 (ii) \(\Rightarrow\) (i)。\(\{e\} = \overline{\{e\}}\)，所以由命题 1 有 (i) \(\Rightarrow\) (ii)。

推论 2. 设 G 是 Hausdorff 拓扑群，A 是 G 的子群。A 可解（分别地，幂零、交换）的充要条件是 \(\overline{\mathbf{A}}\) 具有相同性质。

这立即由命题 1 得到。

命题 2. 设 \(\mathbf{G}\) 是拓扑群，\(\mathbf{A}\) 和 \(\mathbf{B}\) 是 \(\mathbf{G}\) 的子群。若 \(\mathbf{A}\) 连通，则 (A, B) 连通。

固定 B 中的 y。由所有 \(M_{y}\) 构成的 \((x,y)\) 中的 \(x \in A\) 的集合是连通的（因为从 A 到 G 的映射 \(x \mapsto (x,y)\) 是连续的）。由于 \(e \in M_{y}\)，所以所有 \(M_{y}\) 的 \(y \in B\) 的并集 R 是连通的。但 (A, B) 是 G 中由 R 生成的子群，故得命题。

